% Chapter 2: safety.
\chapter{Safety}\label{ch:safety}
\index{safety}

Read this page before you use the machine. Post the quick-reference card
(Appendix~\ref{app:quickref}) beside the panel.

\begin{warning}[WARNING -- MOVING PARTS]
The roll, the motor drive and the measuring wheel can pull in fingers, hair,
loose sleeves, gloves and jewellery. Switch the machine's power off before you
load paper, thread paper, clear a jam, clean the wheel or reach into the
paper path.
\end{warning}

\begin{warning}[WARNING -- THE MOTOR CAN START WITHOUT WARNING]
The motor starts the moment someone pushes \btn{1}, \btn{2} or \btn{4}, or
when a person at a connected PC uses the program's Test mode. During a wind
the motor stops and starts again by itself (the trim bumps,
Section~\ref{sec:trim}). A wind is not over until you hear the beeps and see
the result screen. Always assume the motor can move while power is on.
\end{warning}

\begin{warning}[WARNING -- THE BUTTONS ARE NOT AN EMERGENCY STOP]
\btn{1} and \btn{2} stop a wind through the controller. If the controller
has a fault, the screen is frozen, or the motor does not stop, \textbf{switch
the machine's power off}. The paper keeps moving for a short distance after
the motor stops (the ``coast'', up to a few yards at full speed).
\end{warning}

\textbf{General rules}
\begin{itemize}
  \item Only trained staff should run the machine. Keep children and visitors away.
  \item Manual mode (\btn{4}) runs the motor continuously until you push
        \btn{4} again. Never leave the machine in manual mode unattended.
  \item The machine does \emph{not} detect a jam during the main part of the
        wind (only during the trim). If the paper stops, bunches or tears,
        push \btn{1} or \btn{2} at once, or switch the power off.
  \item If someone is using the PC program on the machine, agree who is in
        charge. While Test mode is on, the panel cannot start a wind, but
        \btn{1} and \btn{2} still stop the motor.
  \item Do not open the control box with the power on. Only a qualified person
        should work on the wiring, the relay or the motor.
  \item Connect only the three wires of the USB serial adapter (TX, RX, GND)
        described in Section~\ref{sec:adapter}. Do not power the machine from
        the adapter.
  \item Keep the work area dry and the cables out of walkways.
  \item Machine power rating and supply: [POWER RATING -- SEE MACHINE LABEL].
\end{itemize}

\textbf{If something goes wrong}: switch the power off, wait until all
motion has stopped, then see Chapter~\ref{ch:trouble}.
